\SourcePage{614}{617}
\section{The unitarity condition in scattering}

The scattering amplitude in an arbitrary field (which need not be central)
obeys certain relations that follow from general physical requirements.

At large distances, the asymptotic form of the wave function for elastic
scattering in an arbitrary field is
\begin{equation}
 \psi\approx e^{ikr\mathbf n\mathbf n'}+\frac1r
 f(\mathbf n,\mathbf n')e^{ikr}.
 \tag{125.1}
\end{equation}
This notation differs from (123.3) in that the scattering amplitude now
depends on the directions of two unit vectors: one along the incident
direction ($\mathbf n$), and the other along the scattered direction
($\mathbf n'$), rather than only on the angle between them.

Any linear combination of functions of the form (125.1), with different
incident directions $\mathbf n$, likewise represents a possible scattering
process. Multiply the functions (125.1) by arbitrary coefficients
$F(\mathbf n)$ and integrate over all directions $\mathbf n$ (with solid-angle
element $d\Omega$). Such a linear combination may then be written as
\begin{equation}
 \int F(\mathbf n)e^{ikr\mathbf n\mathbf n'}d\Omega
 +\frac{e^{ikr}}r\int F(\mathbf n)f(\mathbf n,\mathbf n')d\Omega.
 \tag{125.2}
\end{equation}

\SourcePage{615}{618}
Since $r$ may be arbitrarily large, the factor
$e^{ikr\mathbf n\mathbf n'}$ in the first integral oscillates rapidly as the
direction of the variable vector $\mathbf n$ changes. The integral is
therefore determined mainly by neighborhoods of those values of $\mathbf n$
at which the exponent is extremal ($\mathbf n=\pm\mathbf n'$). Within either
neighborhood, the factor $F(\mathbf n)\approx F(\pm\mathbf n')$ may be taken
outside the integral; integration then gives\footnote{To evaluate the
integral, deform the integration path in the complex $\mu$ plane, where
$\mu=\cos\theta$ ($\theta$ is the angle between $\mathbf n$ and
$\mathbf n'$), so that it bends into the upper half-plane while its endpoints
$\mu=\pm1$ remain fixed. The function $e^{ikr\mu}$ then decays rapidly on
moving away from either endpoint.}
\[
 \frac{2\pi i}{kr}F(-\mathbf n')e^{-ikr}
 -\frac{2\pi i}{kr}F(\mathbf n')e^{ikr}
 +\frac{e^{ikr}}r\int f(\mathbf n,\mathbf n')F(\mathbf n)d\Omega.
\]
Dropping the common factor $2\pi i/k$, rewrite this expression in the compact
operator form
\begin{equation}
 \frac{e^{-ikr}}rF(-\mathbf n')-\frac{e^{ikr}}r\hat S F(\mathbf n'),
 \tag{125.3}
\end{equation}
where
\begin{equation}
 \hat S=1+2ik\hat f,
 \tag{125.4}
\end{equation}
and $\hat f$ is the integral operator
\begin{equation}
 \hat fF(\mathbf n')=\frac1{4\pi}\int
 f(\mathbf n,\mathbf n')F(\mathbf n)d\Omega.
 \tag{125.5}
\end{equation}
The operator $\hat S$ is called the scattering operator (or scattering
\emph{matrix}), or simply the $S$-matrix; it was introduced by
\emph{W.~Heisenberg} in 1943.

The first term in (125.3) is a wave converging on the centre, whereas the
second is a wave diverging from it. Conservation of particle number in
elastic scattering means that the total particle currents in the converging
and diverging waves are equal. Equivalently, these waves must have the same
normalization. The scattering operator must consequently be unitary (see
\S~12), so that
\begin{equation}
 \hat S\hat S^+=1,
 \tag{125.6}
\end{equation}
or, on substituting (125.4) and multiplying out,
\begin{equation}
 \hat f-\hat f^+=2ik\hat f\hat f^+.
 \tag{125.7}
\end{equation}

\SourcePage{616}{619}
Finally, using the definition (125.5), the \emph{unitarity} condition for
scattering takes the form
\begin{equation}
 f(\mathbf n,\mathbf n')-f^*(\mathbf n',\mathbf n)
 =\frac{ik}{2\pi}\int f(\mathbf n,\mathbf n'')
 f^*(\mathbf n',\mathbf n'')d\Omega''.
 \tag{125.8}
\end{equation}
When $\mathbf n=\mathbf n'$, the integral on the right is precisely the total
scattering cross-section
\[
 \sigma=\int|f(\mathbf n,\mathbf n'')|^2d\Omega''.
\]
The difference on the left then reduces to the imaginary part of the
amplitude $f(\mathbf n,\mathbf n)$. We thus obtain the following general
relation between the total elastic-scattering cross-section and the imaginary
part of the forward-scattering amplitude:
\begin{equation}
 \operatorname{Im}f(\mathbf n,\mathbf n)=\frac{k}{4\pi}\sigma
 \tag{125.9}
\end{equation}
(the so-called \emph{optical theorem} for scattering).

Another general property of the scattering amplitude follows from the
requirement of time-reversal symmetry. In quantum mechanics this symmetry
means that, if $\psi$ describes a possible state, the complex-conjugate
function $\psi^*$ also corresponds to a possible state (see \S~18). Hence the
wave function
\[
 -\frac{e^{-ikr}}rF^*(-\mathbf n')
 -\frac{e^{ikr}}r\hat S^*F^*(\mathbf n'),
\]
which is the complex conjugate of (125.3), also describes a possible
scattering process. Introduce a new arbitrary function by
$-\hat S^*F^*(\mathbf n')=\Phi(-\mathbf n')$. From the unitarity of $\hat S$
we have
\[
 F^*(\mathbf n')=-\hat S^{*-1}\Phi(-\mathbf n')=-\widetilde{\hat S}
 \Phi(-\mathbf n');
\]
introducing the coordinate-inversion operator $\hat P$, which reverses the
signs of $\mathbf n$ and $\mathbf n'$, we write
\[
 F^*(-\mathbf n')=-\hat P\widetilde{\hat S}\hat P\Phi(\mathbf n').
\]
The time-reversed wave function therefore becomes
\[
 \frac{e^{-ikr}}r\Phi(-\mathbf n')
 -\frac{e^{ikr}}r\hat P\widetilde{\hat S}\hat P\Phi(\mathbf n').
\]

\SourcePage{617}{620}
In substance, it must coincide with the original wave function (125.3).
Comparison shows that this requires
\begin{equation}
 \hat P\widetilde{\hat S}\hat P=\hat S,
 \tag{125.10}
\end{equation}
after which the two functions differ only in the name given to the arbitrary
function.

The corresponding relation for the scattering amplitude is obtained by
passing from the operator equality (125.10) to its matrix form. Transposition
interchanges the initial and final vectors $\mathbf n$ and $\mathbf n'$, while
inversion reverses their signs. Consequently,
\begin{equation}
 S(\mathbf n,\mathbf n')=S(-\mathbf n',-\mathbf n),
 \tag{125.11}
\end{equation}
or equivalently,
\begin{equation}
 f(\mathbf n,\mathbf n')=f(-\mathbf n',-\mathbf n).
 \tag{125.12}
\end{equation}
This relation (the so-called \emph{reciprocity theorem}) expresses the natural
result that two scattering processes which are time reverses of one another
have equal amplitudes. Time reversal interchanges the initial and final states
and reverses the particles' directions of motion in both.

The general relations just obtained simplify for scattering in a central
field. The amplitude $f(\mathbf n,\mathbf n')$ then depends only on the angle
$\theta$ between $\mathbf n$ and $\mathbf n'$, so (125.12) becomes an
identity. The unitarity condition (125.8), on the other hand, becomes
\begin{equation}
 \operatorname{Im}f(\theta)=\frac{k}{4\pi}
 \int f(\gamma)f^*(\gamma')d\Omega'',
 \tag{125.13}
\end{equation}
where $\gamma$ and $\gamma'$ are the angles made by $\mathbf n$ and
$\mathbf n'$, respectively, with some direction $\mathbf n''$ in space. If
$f(\theta)$ is expanded as in (123.14), the addition theorem for spherical
functions (p.~10) applied to (125.13) gives the following relation for the
partial amplitudes:
\begin{equation}
 \operatorname{Im}f_l=k|f_l|^2.
 \tag{125.14}
\end{equation}
The same formula follows directly from (123.15), according to which
$|2ikf_l+1|^2=1$. For scattering in a central field the optical theorem
(125.9) also follows at once from (123.11) and (123.12).

\SourcePage{618}{621}
Writing (125.14) as $\operatorname{Im}(1/f_l)=-k$, we see that $f_l$ must have
the form
\begin{equation}
 f_l=\frac1{g_l-ik},
 \tag{125.15}
\end{equation}
where $g_l=g_l(k)$ is real; it is related to the phase $\delta_l$ by
\begin{equation}
 g_l=k\cot\delta_l.
 \tag{125.16}
\end{equation}
We shall use this representation of the amplitude repeatedly below.

For scattering in a central field, let us now trace the connection between
the scattering operator introduced above and the quantities used in the
theory of \S~123.

Since orbital angular momentum is conserved in a central field, the
scattering operator commutes with the angular-momentum operator. In other
words, the $S$-matrix is diagonal in the $l$-representation. Because $\hat S$
is unitary, its eigenvalues must have unit modulus; they are therefore of the
form $\exp(2i\delta_l)$, with real $\delta_l$. It is readily seen that these
quantities are the phase shifts of the wave functions, so the eigenvalues of
the $S$-matrix coincide with the quantities $S_l$ introduced in \S~123,
(123.10); correspondingly, the eigenvalues of
$\hat f=(\hat S-1)/(2ik)$ coincide with the partial amplitudes (123.15).
Indeed, choose $P_l(\cos\theta)$ for $F(\mathbf n)$ (then
$F(-\mathbf n)=P_l(-\cos\theta)=(-1)^lP_l(\cos\theta)$). The wave function
(125.3) must coincide with the solution of the Schrödinger equation
represented by a single term of the sum in (123.9). This means precisely that
\[
 \hat SP_l(\cos\theta)=S_lP_l(\cos\theta).
\]

For a plane wave incident along the $z$ axis, the function $F(\mathbf n)$ in
(125.3) is the delta-function $F=4\delta(1-\cos\theta)$, where $\theta$ is the
angle between $\mathbf n$ and the $z$ axis. Here the delta-function is defined
as stated in the note on p.~616, and the coefficient before it is chosen so
that substitution in the right-hand side of (125.5) yields simply
$f(\theta)$ (where now $\theta$ is the angle between $\mathbf n'$ and the
$z$ axis). Expanding the delta-function in the form (124.3),
\begin{equation}
 F=4\delta(1-\cos\theta)=\sum_{l=0}^{\infty}(2l+1)P_l(\cos\theta),
 \tag{125.17}
\end{equation}
and applying $\hat f$ to it gives, as required, the scattering amplitude in
the form (123.14).

\SourcePage{619}{622}
We conclude with one further observation. Mathematically, the unitarity
condition (125.8) shows that an arbitrarily prescribed function
$f(\mathbf n,\mathbf n')$ cannot necessarily serve as the scattering
amplitude for some field. In particular, an arbitrary $f(\theta)$ need not be
the scattering amplitude for any central field. By (125.13), its real and
imaginary parts must satisfy a definite relation. If
$f(\theta)=|f|e^{i\alpha}$ is written with $|f|$ prescribed for all angles,
(125.13) becomes an integral equation from which the unknown phase
$\alpha(\theta)$ can in principle be found. In other words, knowledge of the
scattering cross-section (the square $|f|^2$) at every angle permits, in
principle, reconstruction of the amplitude as well. This reconstruction is
not wholly unique, however: it determines the amplitude only up to the
replacement
\begin{equation}
 f(\theta)\to-f^*(\theta),
 \tag{125.18}
\end{equation}
which leaves (125.13) invariant and of course does not change the
cross-section $|f|^2$ (the transformation (125.18) is equivalent to reversing
the signs of all the phases $\delta_l$ in (123.11) simultaneously). The
ambiguity is removed, however, if the scattering amplitude is considered as a
function not only of angle but also of energy. We shall see below (\S~128,
129) that the analytic properties of the amplitude as a function of energy
are not invariant under (125.18).
